\section{My role}

At 02:14 on 2026-03-31, after our second cancelled field day, I pushed commit \texttt{7eb7cc8}---a fix for a double-scaling bug that had been silently offsetting my GPS target estimates by over 100\,m. I'd been on the project sixteen weeks at that point, and the commit was my 611th. What stopped me wasn't the bug. It was the realisation that \texttt{passive\_watch.py}, the file I was patching, wasn't in my original remit. Neither were the 127 unit tests, the 2508-line interactive simulator, the refactor from 1411 to 754 lines, or the D6 bibliography I'd rebuild two days later. I had been allocated a narrow CV role at the wk12 kickoff: train a YOLOv8n dummy detector and hand it to the flight team.

\textbf{1a.} Everything beyond that core---the simulation framework, the progressive test ladder (bench $\to$ passive $\to$ autonomous, numbered \texttt{0a} through \texttt{4}), the geofence logic, the blueprint system, the formal clarification letter to the course organiser, even Demetro's FDR speaking scripts---I took on implicitly, one commit at a time, whenever I saw something standing between the team and a working demo. My roles didn't change in a single moment; they drifted, silently, across roughly eight weeks during which three cancelled flight days left us with zero hardware to iterate against. Drawing on Lencioni's (2002) five dysfunctions framework, which we studied in our Professional Practice sessions, I now recognise that this drift was partly an \emph{absence-of-trust} pattern made behavioural: building infrastructure alone was a way of feeling productive, and it protected me from the harder skill of slowing down so a teammate could catch up. Tuckman (1965) would say we never adequately Formed---we jumped straight into task division at wk12 and hit Storming at wk17 without having built the interpersonal trust that makes conflict productive. My Ukrainian training told me one person delivers; this project told me that instinct is a failure mode dressed as a strength.

\textbf{1b.} The contribution I'm most proud of sits inside that uncomfortable truth. It's not the retrained model at mAP50\,=\,0.995, not the 820-plus commits, not the D6 Goldmine climb from 74 to 85.2 across six scoring cycles. It's the fact that on the first cancelled flight day (2026-03-11), when weather killed our only window, I didn't go home. I wrote \texttt{FIELD\_QUICK\_REF.md} for no-internet field use, calibrated the focal length to 5.46\,mm against a tape measure at 1\,m, ran the checkerboard lens calibration, and closed ten commits before dark. That day crystallised my one-line thesis about this project: \emph{the simulation framework we built ourselves is the only reason this project produced any verifiable engineering output at all}. What I'm proud of isn't the output---it's the judgement call that a missing drone must not become a missing report. Had I treated the field days as the unit of work, the team would have delivered zero flights \emph{and} zero verifiable data. Instead we delivered zero flights and a working state machine, a 2508-line simulator with GPS drift and camera shake, 146 pytest tests, and a web ground station a pilot could operate unaided. Had I acted on that ``build a simulator first'' instinct a week earlier, we'd have saved roughly five days of idle waiting. I trust the escape route too late---a pattern I now recognise in myself across work and life.

\textbf{1c.} What I'll focus on differently is the distinction between \emph{capability} and \emph{enablement}. I leave this project believing that engineering leadership is less about being the one who fixes the geofence sign-convention bug at midnight (reflection-in-action; Sch\"on, 1983), and more about making my fix legible enough that the teammate who owns the module could have found it without me. The evidence that convinced me was the \texttt{passive\_watch\_2.py} handoff to Robin (commit \texttt{a5b1ab7}): I built the \texttt{--robin} flag and the \texttt{cv\_mode} HTTP polling interface so his state machine could consume my vision output without touching my code---and Robin integrated for two days without asking a single question. That was the first teamwork move I made that scaled without my presence, and it happened in week~20, not week~12. In my robotics internship starting June I'll schedule two hours per week of explicit pair-programming, and the falsifiable signal I'll watch for is whether a teammate can complete a handoff task independently within 48 hours of receiving it. If they can't, my documentation has failed, not their attention.
